\chapter{Fan-Out: Parallel Budget Relays and Cross-Branch Disagreement}
\label{ch:fanout}

\begin{quotation}
\noindent\textit{Synopsis.} The framework is extended here from linear
chains to branching structures, where a single source generates multiple
downstream renderings optimized for different purposes. Branch agreement
and silent disagreement are introduced to describe situations in which
branches appear compatible while relying on incompatible assumptions about
discarded information. Confidence and quality estimates are shown to be
renderings in their own right, subject to the same limitations as the
artifacts they describe, which motivates the chapter's closing move: a
trust boundary as the practical stopping point that prevents an infinite
regress in evaluating confidence and provenance systems.
\end{quotation}

\section{From Chains to Graphs}

\begin{definition}[Fan-Out Relay]
\label{def:fanout-relay}
A fan-out relay is a directed graph of renderings with a single source
$A_0$ and two or more independent successor branches
$A_0 \xrightarrow{B^{(1)}} A^{(1)}$, $A_0 \xrightarrow{B^{(2)}} A^{(2)}$,
\ldots, each governed by its own local budget triple and each satisfying
\cref{def:budget-relay}'s independence condition with respect to every
other branch: branch $k$'s transition is selected without reference to
any other branch's budget or ledger.
\end{definition}

\cref{def:fanout-relay} generalizes \cref{def:substrate-chain} directly:
audio rendered in parallel into a transcript, an acoustic-feature
analysis, and an emotion analysis is the paradigm case, and each branch
may itself continue as a linear substrate chain of arbitrary length
beyond the point of divergence. What is new here is not the branching
structure by itself, but what it makes possible: two renderings that
both trace back to the same source, and that a downstream party might
reasonably expect to agree with each other, can instead diverge in ways
that neither branch's own local admissibility (\cref{def:locally-admissible})
does anything to prevent.

\section{Branch Agreement and Silent Disagreement}

\begin{definition}[Branch Agreement]
\label{def:branch-agreement}
Two renderings $A^{(i)}, A^{(j)}$ derived from a common ancestor are in
agreement on a distinction $d$ if neither branch's discard set
contradicts the other's claims about $d$.
\end{definition}

\begin{definition}[Silent Disagreement]
\label{def:silent-disagreement}
Branches disagree silently when each branch's confidence or quality
profile $Q$ fails to flag that another branch depended on a distinction
that the first branch discarded or rendered unstable.
\end{definition}

The definitions above are easy to accept in the abstract and easy to
under-examine in practice, so this section works a single construction
in full, in the same style as Appendix~\ref{app:proof}'s canonical
witness, rather than resting on a description of the transcript/acoustic
case alone.

\subsection{A Minimal Two-Branch Witness}

\begin{construction}[Two-Branch Witness]
\label{con:two-branch}
Let $A_0$ carry two distinctions: $d$, the resolved identity of an
ambiguous phoneme, with cost $C(d) = 2$; and $q_d$, a meta-distinction
recording whether the source evidence for $d$ is stable or unstable, with
cost $C(q_d) = 1$. Branch $X$ (a transcript) has task $\tau_X$ requiring
exactly $\{d\}$, with budget $\beta_X = 2$. Branch $Y$ (an
acoustic-stability analysis) has task $\tau_Y$ requiring exactly
$\{q_d\}$, with budget $\beta_Y = 1$.
\end{construction}

\begin{proposition}[The Witness Produces Silent Disagreement]
\label{prop:witness-silent-disagreement}
In \cref{con:two-branch}, branch $X$ is locally admissible with respect
to $\tau_X$ and discards $q_d$; branch $Y$ is locally admissible with
respect to $\tau_Y$ and discards $d$; and a consumer with access only to
branch $X$ cannot detect that $q_d$ -- and hence the instability branch
$Y$ would have reported -- was ever at issue.
\end{proposition}

\begin{proof}
For branch $X$: the required distinction $\{d\}$ has cost $R_X = C(d) = 2
= \beta_X$ exactly. The only non-required distinction present, $q_d$, has
cost $C(q_d) = 1 > \beta_X - R_X = 0$. By \cref{lem:exhaustion}, $q_d$ is
dropped under any construction satisfying \cref{def:budget-discard} with
$d$ ranked above $q_d$, so $X$ is locally admissible
(\cref{def:locally-admissible}) and $D_X = \{q_d\}$.

For branch $Y$: symmetric. The required distinction $\{q_d\}$ has cost
$R_Y = C(q_d) = 1 = \beta_Y$ exactly, and the non-required $d$ has cost
$C(d) = 2 > \beta_Y - R_Y = 0$, so by \cref{lem:exhaustion}, $d$ is
dropped under any construction, $Y$ is locally admissible, and $D_Y =
\{d\}$.

By \cref{prop:confidence-renderings}, $X$'s confidence profile $Q_X$ is
itself a rendering, produced from whatever $X$ retains; since $q_d \notin
A^{(X)}$, $Q_X$ has no distinction available from which to construct an
instability flag, and $X$'s output presents $d$'s resolved value without
qualification. A consumer relying on branch $X$ alone therefore receives
no signal --- not a low-confidence flag, not a declared gap, nothing ---
indicating that $q_d$ was ever a live consideration, even though branch
$Y$, derived from the identical source $A_0$, would have reported it as
unstable. This is exactly \cref{def:silent-disagreement}'s condition:
$Q_X$ fails to flag that another branch's finding (instability, via
$q_d$) bears on distinction $d$, which $X$ itself depended on.
\end{proof}

This construction is deliberately built to reuse
Appendix~\ref{app:proof}'s exhaustion lemma rather than argue the point
freshly, and the reuse is informative in its own right: the same
mechanism that forces a construction-independent discard within a single
linear chain (\cref{lem:exhaustion}) forces construction-independent,
\emph{asymmetric} discard across two branches of a fan-out, and it is
precisely this asymmetry --- $X$ drops what $Y$ keeps, and vice versa ---
that produces silent disagreement rather than mere independent loss.

\begin{proposition}[Silent Disagreement Exists in General]
\label{prop:silent-disagreement}
Given two branches produced independently (\cref{def:fanout-relay}) from
a common source, there exist source distinctions whose instability is
reported by one branch and concealed by the other.
\end{proposition}

\begin{proofsketch}
\cref{con:two-branch} together with \cref{prop:witness-silent-disagreement}
constitutes an existence proof: at least one pair of branches, with at
least one such distinction, is exhibited explicitly and shown to satisfy
\cref{def:silent-disagreement} under a construction-independent argument.
This result was labeled a Theorem in an earlier draft; the label has been
corrected to Proposition, consistent with the monograph's
epistemic-labeling discipline, since an existence claim witnessed by one
explicit construction is not of the same weight as
\cref{thm:composition-failure}'s general composition result.
\end{proofsketch}

\section{The Regress Inside the Confidence Profile}

\begin{proposition}[Confidence Profiles Are Renderings]
\label{prop:confidence-renderings}
A confidence or quality profile $Q$ attached to a rendering $A$ is itself
produced by a projection, and therefore has its own discard set
$\discard{Q}$.
\end{proposition}

\begin{proofsketch}
Confidence estimation maps source evidence -- acoustic likelihoods, model
logits, annotator judgments -- into a finite label or score. Any such
finite mapping collapses distinctions in the sense of
\cref{def:rendering}; $Q$ therefore satisfies the definition of a
rendering and inherits a discard set like any other.
\end{proofsketch}

\cref{prop:confidence-renderings} produces a regress: trusting $Q$
requires trusting the process that produced $Q$, which has its own
confidence measure $Q'$, and so on. This regress does not resolve within
the formalism; it is cut off by convention rather than by argument, and
it is worth being explicit about the cost of extending it even one level
further. Verifying $Q$ against a further $Q'$ is itself a repair
decision in the sense of Chapter~\ref{ch:repair-economics}: it consumes
budget that could otherwise be spent on the object-level distinctions
$Q$ was meant to qualify, and \cref{prop:no-global-conserved-quantity}'s
argument applies just as well to a chain of nested confidence measures as
to a chain of substantive renderings. Real systems stop extending this
chain not because the regress terminates on its own but because someone
decides, by convention, where it is no longer worth the expenditure.

\begin{definition}[Trust Boundary]
\label{def:trust-boundary}
A designated layer -- for instance, a fixed, versioned confidence model
-- beyond which a relay does not recursively question its own
loss-accounting. The trust boundary is a declared convention, not a
theorem; \cref{prop:confidence-renderings}'s regress is halted by
stipulation at this layer, and an auditor of the relay is, in effect,
trusting the boundary itself rather than anything beyond it.
\end{definition}

\begin{definition}[Boundary-Relative Sufficiency]
\label{def:boundary-sufficiency}
A rendering is boundary-sufficient for a task if it is sufficient
(\cref{def:sufficiency}) conditional on every confidence and quality
profile at or beyond the declared trust boundary
(\cref{def:trust-boundary}) being accepted without further verification.
\end{definition}

\cref{def:boundary-sufficiency} names what practitioners actually rely
on, as distinct from the stronger and generally unattainable notion of
unconditional sufficiency: a system is almost never verified all the way
down to its source distinctions, and boundary-sufficiency makes explicit
exactly how much trust that verification is quietly resting on.
